\section{Verifiche numeriche: modello 2}\label{app:num2}

Gate di prova $U=(r_1\otimes r_2)\,e^{-i\tau(XX+YY+JZZ)}(r_2^\dagger\otimes r_1)$, con $\tau=0.5$, $J=0.3$ e rotazioni a un qubit fissate. L'autovalore subdominante del canale a una replica è $\lambda=0.4622$. Per i qutrit si usa un gate fisso casuale, con $\lambda=0.394$.

\paragraph{Ricorsione esatta contro Monte Carlo} (media su $V_s$ Haar; $A=B=Z$ per $k=2$, $A=Z+\tfrac12$, $B=X+\tfrac12$ per $k=3$).
\begin{center}\small
\begin{tabular}{@{}lccc@{}}
\toprule
caso & $t$ & Monte Carlo & esatto\\
\midrule
$k=2$, $D=2$, $p=0$ & 2 / 3 & $-0.0788(32)$ / $-0.0526(49)$ & $-0.0825$ / $-0.0483$\\
$k=2$, $D=3$, $p=0$ & 2 / 3 & $-0.0689(21)$ / $-0.0355(31)$ & $-0.0689$ / $-0.0373$\\
$k=2$, $D=2$, rumore su E $p=0.3$ & 3 & $-0.0460(30)$ & $-0.0452$\\
$k=2$, $D=2$, rumore su C $p=0.3$ & 3 & $-0.0260(17)$ & $-0.0247$\\
$k=3$, $D=2$, $p=0$ & 2 / 3 & $1.0828(93)$ / $1.0390(90)$ & $1.0865$ / $1.0310$\\
$k=3$, $D=2$ / $3$, $p=0.4$ & 3 & $0.6394(51)$ / $0.5993(30)$ & $0.6450$ / $0.5960$\\
\bottomrule
\end{tabular}
\end{center}
Il limite a due settori coincide con la ricorsione esatta a $D=10^3$ e $10^5$, a 6 cifre (per esempio $C_2(5)=-0.001840$ senza rumore e $-0.001189$ con $p=0.2$). Il codice generale in $k$ riproduce quello dedicato a $k=2$.

\paragraph{Valori a tempi lunghi} ($t=60$, $A=Z+a$, $B=X+b$, $k=2$). Valore libero $a^2+b^2+a^2b^2$, valore monotono $a^2(1+b^2)$.
\begin{center}\small
\begin{tabular}{@{}ccccccc@{}}
\toprule
$(a,b)$ & $p=0$ & $0.02$ & $0.1$ & $0.3$ & libero & monotono\\
\midrule
$(0.5,0)$ & 0.2500 & 0.2500 & 0.2500 & 0.2500 & 0.2500 & 0.2500\\
$(0,0.5)$ & 0.2500 & 0.0242 & 0.0000 & 0.0000 & 0.2500 & 0\\
$(0.5,0.7)$ & 0.8625 & 0.4199 & 0.3725 & 0.3725 & 0.8625 & 0.3725\\
$(1,1)$ & 3.0000 & 2.0967 & 2.0000 & 2.0000 & 3.0000 & 2.0000\\
\bottomrule
\end{tabular}
\end{center}
Con rumore su C ($p=0.3$ e $0.8$) o su A ($p=0.3$) si ottiene lo stesso valore monotono, $0.3125$ per $(a,b)=(0.5,0.5)$.

A $k=3$ (qubit, $a=b=\tfrac12$): valore libero $1.140625$, monotono $0.203125$. Si trova $1.140625$ senza rumore. Con $p=0.1$ il valore tende a $0.2031$ con tasso $0.809\approx(1-p)^2$; con $p=0.3$, a $t=40$, si ha $0.203125$. Con i qutrit, $A=\mathrm{diag}(1,1,-2)$ e $B=\mathrm{diag}(1,0,0)$: il valore libero $\kappa_3(A)\varphi(B)^3=-0.074074$ è riprodotto senza rumore, e con $p=0.1$ il segnale decade con tasso $0.7290=(1-p)^3$. A $k=4$ (qubit, $A=Z$, $B=X+\tfrac12$): $C_4(14)=C_4(20)=0.562500$ senza rumore, che è il valore libero $2b^2+b^4$. Con $p=0.1$ il tasso è $0.6561=(1-p)^4$.

\paragraph{Tassi di rilassamento} ($k=2$, $A=Z$, $B=X+\tfrac12$, tasso stimato su $t\in[14,24]$).
\begin{center}\small
\begin{tabular}{@{}lccccccc@{}}
\toprule
$p$ & 0.1 & 0.3 & 0.45 & 0.5 & 0.55 & 0.6 & 0.7\\
\midrule
tasso & 0.8100 & 0.4900 & 0.3028 & 0.2525 & 0.2207 & 0.2142 & 0.2140\\
$\max(u,\lambda^2)$ & 0.8100 & 0.4900 & 0.3025 & 0.2500 & 0.2136 & 0.2136 & 0.2136\\
\bottomrule
\end{tabular}
\end{center}
Per $A=Z+\tfrac12$ il tasso misurato è invece $\max(u,\lambda)$: $0.4866$ a $p=0.3$ ($u=0.49$), e $0.459$--$0.460$ per $p\ge0.45$ ($\lambda=0.462$).
